% Theme: Textbook  --  best for definitions, theorems and short proofs
% Build:  ./build.sh example-textbook.tex   (or compile on Overleaf)
\documentclass[12pt]{beamer}
\usetheme{igtextbook}

\ighandle{@quantum_danesh}
\igseries{Corrección de errores}
\igissue{04}

\begin{document}

\coverblurb{Qué es un código, cuándo un error es corregible y cómo
comprobarlo con una ecuación.}
\coverslide{Capítulo 4}
  {Las condiciones de Knill--Laflamme}
  {¿Cuándo se puede deshacer un error?}

\begin{frame}{Códigos y errores}{4.1}
\begin{defbox}[Definición 4.1]
Un \key{código cuántico} es un subespacio $\mathcal{C}\subseteq\mathcal{H}$
con proyector $P$. Un conjunto de errores $\{E_a\}$ es \key{corregible} si
existe un canal $\mathcal{R}$ tal que
\[
  \mathcal{R}\Big(\sum_a E_a\,\rho\,E_a^\dagger\Big) = \rho
\]
para todo estado $\rho$ con soporte en $\mathcal{C}$.
\end{defbox}
La pregunta es cómo saberlo sin construir $\mathcal{R}$.
\end{frame}

\begin{frame}{El teorema}{4.2}
\begin{thmbox}[Teorema 4.2 (Knill--Laflamme)]
$\{E_a\}$ es corregible en $\mathcal{C}$ si y solo si
\[
  P\,E_a^\dagger E_b\,P = \lambda_{ab}\,P \qquad \text{para todo } a,b,
\]
con $(\lambda_{ab})$ una matriz hermítica.
\end{thmbox}
Equivalentemente, para una base ortonormal $\{\ket{i}\}$ de $\mathcal{C}$:
\keyeq{\mel{i}{E_a^\dagger E_b}{j} = \lambda_{ab}\,\delta_{ij}}
\end{frame}

\begin{frame}{Cómo leerlo}{4.2}
\[
  \mel{i}{E_a^\dagger E_b}{j}
  = \explain<2->{\mhl<2->{\lambda_{ab}}}{igual para todo $i$}\;
    \explain<3->{\mhl*<3->{\delta_{ij}}}{no mezcla palabras}
\]
\begin{itemize}
  \item<2-> $\lambda_{ab}$ no depende de $i$: el entorno no aprende
        \hl{nada} sobre el estado lógico.
  \item<3-> $\delta_{ij}$: los errores mantienen ortogonales las palabras
        código, así que se pueden distinguir.
\end{itemize}
\end{frame}

\begin{frame}{Ejemplo: el código de repetición}{4.3}
Sea $\mathcal{C} = \operatorname{span}\{\ket{000},\ket{111}\}$ con errores
$\{\id, X_1, X_2, X_3\}$.
\begin{igproof}[Comprobación]
Si $a\neq b$, $E_a^\dagger E_b$ es un producto de una o dos $X$ y saca
$\ket{000}$ y $\ket{111}$ del código, así que
$\mel{i}{E_a^\dagger E_b}{j} = 0$. Si $a = b$, $E_a^\dagger E_a = \id$.
\end{igproof}
\keyeq{\lambda_{ab} = \delta_{ab}}
Corrige cualquier error de un solo bit.
\end{frame}

\begin{frame}{Códigos degenerados}{4.4}
\begin{notebox}[Observación 4.4]
$\lambda$ no tiene por qué ser invertible. Si no lo es, el código es
\key{degenerado}: errores distintos actúan igual sobre el código y no hace
falta distinguirlos.
\end{notebox}
El código de Shor de 9 qubits es el ejemplo clásico: $Z_1$ y $Z_2$ hacen
lo mismo sobre sus palabras código.
\vfill
{\footnotesize\color{igMuted}E.~Knill, R.~Laflamme, Phys.\ Rev.\ A \textbf{55}, 900 (1997).\par}
\end{frame}

\closingslide{Una ecuación para decidir}
  {Próximo capítulo: los códigos estabilizadores, donde estas condiciones
   se comprueban contando conmutadores.}

\end{document}
